% Part: second-order-logic
% Chapter: metatheory
% Section: introduction

\documentclass[../../../include/open-logic-section]{subfiles}

\begin{document}

\olfileid{sol}{met}{int}

\olsection{પ્રસ્તાવના}

પ્રથમ-ક્રમ તર્કશાસ્ત્રમાં કેટલાક સારા ગુણધર્મો છે. આપણે જાણીએ છીએ કે
તે નિર્ણેય નથી, પરંતુ ઓછામાં ઓછું સ્વયંસિદ્ધીકરણીય છે. એટલે કે
પ્રથમ-ક્રમ તર્કશાસ્ત્ર માટે યથાર્થ અને પૂર્ણ સાબિતી-તંત્રો છે: તેઓ
એવો !!a{derivability} સંબંધ~$\Proves$ આપે છે કે !!{sentence}sના કોઈપણ
ગણ~$\Gamma$ અને !!{sentence}~$!A$ માટે $\Gamma \Entails !A$ ત્યારે અને માત્ર
ત્યારે જ જ્યારે $\Gamma \Proves !A$. ખાસ કરીને તેનો અર્થ એ થાય છે કે
પ્રથમ-ક્રમ તર્કશાસ્ત્રનાં પ્રમાણભૂત વાક્યો !!{computably
  enumerable} છે. એવું સંગણનીય વિધેય~$f\colon \Nat \to
\Sent[L]$ છે કે~$f$નાં મૂલ્યો~$\Lang{L}$નાં બધાં અને માત્ર પ્રમાણભૂત
!!{sentence}s છે. આવું બને છે કારણ કે !!{derivation}sની પરિગણના
કરી શકાય છે અને એક !!{sentence} !!{derive} કરતી નિષ્પત્તિઓને પછી
તે !!{sentence} સાથે જોડી શકાય છે. દ્વિતીય-ક્રમ તર્કશાસ્ત્ર
પ્રથમ-ક્રમ તર્કશાસ્ત્ર કરતાં વધુ અભિવ્યક્તિશીલ છે; તેથી તેનાં
પ્રમાણભૂત વાક્યોનું નિરૂપણ સામાન્ય રીતે વધુ જટિલ છે. હકીકતમાં આપણે
બતાવીશું કે દ્વિતીય-ક્રમ તર્કશાસ્ત્ર માત્ર અનિર્ણેય જ નથી, પણ તેના
પ્રમાણભૂત વાક્યો !!{computably
  enumerable} પણ નથી. તેથી દ્વિતીય-ક્રમ તર્કશાસ્ત્ર માટે યથાર્થ અને
પૂર્ણ સાબિતી-તંત્ર હોઈ શકે નહિ (જોકે યથાર્થ પરંતુ અપૂર્ણ સાબિતી-તંત્રો
ઉપલબ્ધ છે અને તે સંશોધનના મહત્વના વિષયો પણ છે).

પ્રથમ-ક્રમ તર્કશાસ્ત્રના બીજા બે ગુણધર્મો પણ છે: તે સઘન છે (!!{sentence}sના
ગણ~$\Gamma$નો દરેક સાન્ત ઉપગણ સંતોષ્ય હોય તો~$\Gamma$ પોતે સંતોષ્ય
છે) અને તેના માટે લેવેનહાઇમ--સ્કોલેમ પ્રમેય સાચું છે (જો~$\Gamma$નું
અનંત નિદર્શ હોય તો તેનું !!a{denumerable} નિદર્શ પણ હોય છે).
દ્વિતીય-ક્રમ તર્કશાસ્ત્રમાં આ બંને પરિણામો નિષ્ફળ જાય છે. અહીં પણ
કારણ એ જ છે કે દ્વિતીય-ક્રમ તર્કશાસ્ત્ર !!{domain}sના કદ વિશે એવી
હકીકતો વ્યક્ત કરી શકે છે જે પ્રથમ-ક્રમ તર્કશાસ્ત્ર કરી શકતું નથી.

\end{document}
